Shelfmark 149, scan 126.

\scanpage{126}
whose residue field (contained in that of $V$) will be \deleted{algebraic} finite over $K$ if that of $V$ is (hence $V_1\ne L_1$) --- hence $V_1$ \deleted{$\longrightarrow$} \inserted{corresponds to} a ``place'' of the function field in one variable $L_1$ over $K$. Now, by essentially obvious functorialities, one sees that the homomorphism $\mathcal E_L^{\overline L}\longrightarrow\mathcal E_{L_1}^{\overline{L_1}}$ sends $N_V$ into $N_{V_1}$, $I_V$ into $I_{V_1}$ --- hence the image of $N_V$ normalizes $I_{V_1}$, hence (if one has a \inserted{germ of a} section of $N_V$ over $\mathcal E_K$) the resulting germ of a section of $\mathcal E_{L_1}$ over $\mathcal E_K$ normalizes $I_{V_1}$ --- and its restriction to $\mathcal E_K^0$ centralizes $I_{V_1}$ ---
\margin{one must check that the field of fractions of $V_1$ is indeed $L_1$ --- this is a general fact\ldots}

\uline{Conjecture} Let $K,L$ be finitely generated extensions of $\mathbf Q$, $K\subset L$. Then

a) Every section of $\mathcal E_L$ over $\mathcal E_K$ (or a germ of such, it amounts to the $=$\ldots) normalizes an $\unsure{I_V|N_V}$ associated to a valuation ring \inserted{$V$} of $L$ containing $K$, with residue field algebraic \inserted{over $K$} \deleted{\unsure{,}} and \deleted{\unsure{this}} $V$ is uniquely
\margin{125}
